% =============================================================================
\section{Method of Solution}
\label{sec:method}
% =============================================================================

% -----------------------------------------------------------------------------
\subsection{DAE Structure}
\label{sec:method:structure}
% -----------------------------------------------------------------------------

When all tanks operate in Configuration~A or~C, \cref{eq:mass-balance,eq:energy-balance,eq:solid-temp}
form a pure system of ODEs:
\begin{equation}
  \dot{\vect{x}} = \vect{f}(\vect{x}, t).
\end{equation}
When one or more tanks enter Configuration~B, the discharge edge flow rate is
no longer a prescribed function of time but is instead constrained by
\cref{eq:config-b-clip}.  The system becomes a \emph{semi-explicit index-1 DAE}:
\begin{align}
  \dot{\vect{x}} &= \vect{f}(\vect{x}, \vect{z}, t), \label{eq:dae-ode}\\
  \vect{0}       &= \vect{g}(\vect{x}, \vect{z}, t), \label{eq:dae-alg}
\end{align}
where $\vect{z}$ collects the algebraically determined discharge rates of all
Config-B tanks.  The constraint \cref{eq:dae-alg} is \cref{eq:config-b-clip}
for each active Config-B tank.

Because the constraint is linear in $\vect{z}$ (the clip operation is
monotone), \cref{eq:dae-alg} is solved \emph{explicitly} at the beginning of
each RHS evaluation rather than via an inner Newton iteration — reducing the
DAE to a sequence of modified ODEs.

% -----------------------------------------------------------------------------
\subsection{Per-Step Algebraic Solve for Configuration~B}
\label{sec:method:configB-solve}
% -----------------------------------------------------------------------------

\Cref{alg:rhs} summarises the right-hand-side evaluation procedure.

\begin{figure}[ht]
\begin{tikzpicture}[
    box/.style    = {draw, rounded corners=3pt, fill=blue!5, minimum width=8cm,
                     minimum height=0.7cm, align=left, font=\small, inner sep=6pt},
    diamond/.style= {draw, diamond, fill=yellow!15, minimum width=2.5cm,
                     minimum height=0.7cm, align=center, font=\small, inner sep=4pt},
    arr/.style    = {-{Latex[length=5pt]}, thick},
    node distance = 0.35cm,
  ]
  \node[box] (s1)  {\textbf{1.} Reconstruct $\vect{x}_i = (m_i, T_i, T_{s,i})$ from state vector $\vect{y}$.};
  \node[box] (s2) [below=of s1]
                   {\textbf{2.} Compute all coupling edge flows $\mdot_e^{\text{coup}}$.};
  \node[box] (s3) [below=of s2]
                   {\textbf{3.} For each tank~$i$: determine configuration (A/B/C).};
  \node[box] (s4) [below=of s3]
                   {\textbf{4.\ Config~B tanks:} evaluate $\mdot_\text{disc}^{\max}$ via
                    \cref{eq:config-b-solution}; clip $\mdot_e^{\text{disc}}$.};
  \node[box] (s5) [below=of s4]
                   {\textbf{5.\ Config~C tanks:} compute vent edge flow $\mdot_e^{\text{vent}}$
                    via \cref{eq:vent-rate}.};
  \node[box] (s6) [below=of s5]
                   {\textbf{6.} Compute thermal model: $\dot{Q}_{s,i}$, $\dot{T}_{s,i}$.};
  \node[box] (s7) [below=of s6]
                   {\textbf{7.} Assemble $\dot{m}_i$, $\dot{T}_i$, $\dot{T}_{s,i}$
                    from \cref{eq:mass-balance}, \cref{eq:energy-balance}, and \cref{eq:solid-temp}.};
  \node[box] (s8) [below=of s7]
                   {\textbf{8.} Return $\dot{\vect{y}}$.};

  \foreach \a/\b in {s1/s2, s2/s3, s3/s4, s4/s5, s5/s6, s6/s7, s7/s8}
    \draw[arr] (\a.south) -- (\b.north);
\end{tikzpicture}
\caption{RHS evaluation procedure for one ODE time step.  Steps 4 and 5 handle
the algebraic parts of the DAE; all remaining steps are explicit ODE evaluations.}
\label{alg:rhs}
\end{figure}

Steps 2--5 together constitute the algebraic solve.  The coupling flows (step~2)
are computed using the current state and are treated as explicit inputs to the
remainder of the RHS; no fixed-point iteration is required for the coupling
because the valve models are explicit functions of the state.

% -----------------------------------------------------------------------------
\subsection{Time Integration}
\label{sec:method:integration}
% -----------------------------------------------------------------------------

The modified ODE system (after Config-B discharge throttling has been applied)
is integrated using the \textsc{lsoda} solver from \texttt{scipy.integrate},
accessed through the \texttt{solve\_ivp} interface.

\textsc{lsoda} was chosen for the following reasons:

\begin{itemize}[nosep, left=0pt]
  \item It switches automatically between Adams (non-stiff) and BDF (stiff)
        formulæ, which is important because stiffness changes throughout the
        mission as configurations switch and coupling valves open or close.
  \item Relative and absolute tolerances are set by the user (\texttt{rtol},
        \texttt{atol}) and are tightened by default to \num{1e-4} and
        \num{1e-7} respectively to control mass-conservation drift.
  \item It produces dense output, enabling post-processing at arbitrary time
        points without re-running the simulation.
\end{itemize}

The integrator step size is fully adaptive; no explicit step-size control is
needed from user code.

% -----------------------------------------------------------------------------
\subsection{Stopping Criteria}
\label{sec:method:stopping}
% -----------------------------------------------------------------------------

Integration terminates when any of the following events is detected:

\begin{center}
\begin{tabularx}{\linewidth}{@{} l l X @{}}
\toprule
\textbf{Criterion} & \textbf{Variable} & \textbf{Description} \\
\midrule
Minimum density  & $\rho_i \le \rho_{\min}$ &
  The tank is effectively empty for the given mission.
  Default: $\rho_{\min} = \SI{5.8}{\kilogram\per\cubic\metre}$ (gaseous H$_2$). \\[4pt]
Target density   & $\rho_i \ge \rho_{\text{target}}$ &
  Used in refuel missions.  Integration stops when the tank has been filled
  to the desired density. \\[4pt]
Duration         & $t \ge t_{\text{mission}}$ &
  Hard time limit derived from the mission profile duration. \\
\bottomrule
\end{tabularx}
\end{center}

All criteria are checked at each RHS evaluation.  When a criterion is met, a
flag is set and the RHS returns a near-zero vector, allowing the solver to
cleanly reach the termination point without an abrupt step.

% -----------------------------------------------------------------------------
\subsection{Configuration Switching and Hysteresis}
\label{sec:method:hysteresis}
% -----------------------------------------------------------------------------

Configuration detection must be numerically robust to avoid chattering near
transition boundaries.  The following hysteresis scheme is used for Config~B:

\begin{align}
  \text{A} \to \text{B} &: \quad P_i \le P_{\min,i}(1 - \delta), \\
  \text{B} \to \text{A} &: \quad P_i >  P_{\min,i}(1 + \delta),
\end{align}
with $\delta = 0.01$ (1\%).  The current configuration is stored as instance
state across RHS calls so that the hysteresis condition can reference the
previous configuration.

Config~C (venting) has no hysteresis: it activates immediately when
$P_i \ge P_{\text{vent},i}$ and deactivates when the vent flow brings the
pressure back below the threshold.
